The proof is in Gauge Theory Gravity, with forms read as multivector fields on flat spacetime: a $k$-form is a field of background $k$-vectors, $\alpha = \tfrac{1}{k!}\alpha_{\mu\ldots}\,e^{\mu}\wedge\cdots$, the exterior derivative is $\nabla\wedge$, and the wedge of forms is the outer product of their background parts. An internal gauge field is a one-form $\mathcal{A} = e^{\mu}A_{\mu}$ whose values $A_{\mu}$ lie in the internal Lie algebra, realised as bivectors of an internal Clifford algebra (a compact Lie algebra is a Lie subalgebra of the bivectors of some $\Cl(n,0)$, acting on its own vectors), multiplied by the internal geometric product, which commutes with the background factor; the trace form is the scalar part $\langle\cdot\rangle_{\mathrm{int}}$ of that product, cyclic, and proportional to the Killing form on each simple factor. The field strength is $\mathcal{F} = \nabla\wedge\mathcal{A} + \mathcal{A}\wedge\mathcal{A}$, a background bivector with values in the Lie algebra.

(1) Internal invariance of $I$ is structural: $I$ is a vector of the tractor factor $\Cl(2,4)$, the internal groups act on the internal factor, and $\mathcal{F}$ is a background bivector with internal Lie-algebra values, with no tractor slot for $\epsilon$ or $I$ to contract. Topologicality: the Bianchi identity $D\wedge\mathcal{F} := \nabla\wedge\mathcal{F} + \mathcal{A}\wedge\mathcal{F} - \mathcal{F}\wedge\mathcal{A} = 0$ holds identically ($\nabla\wedge\nabla = 0$ and associativity), so $\nabla\wedge\langle\mathcal{F}\wedge\mathcal{F}\rangle_{\mathrm{int}} = 2\langle D\wedge\mathcal{F}\wedge\mathcal{F}\rangle_{\mathrm{int}} = 0$, using that background bivectors commute under the outer product and the scalar part is cyclic; on a contractible region $\langle\mathcal{F}\wedge\mathcal{F}\rangle_{\mathrm{int}} = \nabla\wedge\mathcal{C}$ with the Chern-Simons trivector $\mathcal{C} = \langle\mathcal{A}\wedge\nabla\wedge\mathcal{A} + \tfrac23\,\mathcal{A}\wedge\mathcal{A}\wedge\mathcal{A}\rangle_{\mathrm{int}}$; and the variation is $\delta\int\langle\mathcal{F}\wedge\mathcal{F}\rangle_{\mathrm{int}} = 2\int\langle D\wedge\delta\mathcal{A}\wedge\mathcal{F}\rangle_{\mathrm{int}} = 2\int\nabla\wedge\langle\delta\mathcal{A}\wedge\mathcal{F}\rangle_{\mathrm{int}} + 2\int\langle\delta\mathcal{A}\wedge D\wedge\mathcal{F}\rangle_{\mathrm{int}}$, a boundary term plus zero: no field equations. The background metric is never used.

(2) Let $(\hb, \Omega)$ be a gravitational field with derived metric $g(a,b) = \hb^{-1}(a)\cdot\hb^{-1}(b)$. The physical field strength of a two-form $\mathcal{F}$ is $\hb^{T}(\mathcal{F})$ (on physical vectors $A = \hb^{-1}(a)$, $B = \hb^{-1}(b)$: $(a\wedge b)\cdot\mathcal{F} = (A\wedge B)\cdot\hb^{T}(\mathcal{F})$), and the Hodge dual of a physical bivector is multiplication by the pseudoscalar $I$. Read back on the background with $\hb^{-T}(X) = I\,\hb(I^{-1}X)/\det\hb$, and normalised to the background orientation as the component formula is,
\[
{*}_{g}\mathcal{F} \;=\; \frac{1}{|\det\hb|}\; I\,\mathcal{G}(\mathcal{F}), \qquad \mathcal{G} = \hb\,\hb^{T},
\]
where $\mathcal{G}$ is the cometric map ($g^{\mu\nu} = e^{\mu}\cdot\mathcal{G}(e^{\nu})$), acting on bivectors as an outermorphism, and $|\det\hb|^{-1} = \sqrt{|g|}$. A change of scale $g \mapsto \lambda^{-2}g$ is $\hb \mapsto \lambda\hb$: $|\det\hb|^{-1}$ gains $\lambda^{-4}$ and $\mathcal{G}$ on bivectors gains $\lambda^{4}$, so ${*}_{g}$ is unchanged. A rotation-gauge change $\hb(a) \mapsto \hb(\tilde R aR)$ leaves $\mathcal{G}$ and $\det\hb$ unchanged. Since $\mathcal{G}$ determines $g$ ($g(a, b) = a\cdot\mathcal{G}^{-1}(b)$), ${*}_{g}$ depends on $[g]$ only, and not on the scale that the frozen point fixes; it does depend on $[g]$ (the check exhibits two non-conformal fields with different stars), and ${*}_{g}^{2} = -1$, from $\mathcal{G}(I\,\mathcal{G}(\mathcal{F})) = \det(\mathcal{G})\,I\,\mathcal{F}$ and $\det\mathcal{G} = (\det\hb)^{2}$. Gravity's action, by contrast, contains no inverse frame (\claimref{dark-energy.frozen-point-makes-stelle-west-vector}(1)).

(3) is (2) read at a corank-one wall. There the cometric acquires a one-dimensional radical and the metric a pole (Fact \claimref{dark-energy.imported-walls-sphere-theorem}): along a family approaching the wall, $\mathcal{G}_{s} \to \mathcal{G}_{0}$ of rank three and $\det\hb_{s} \to 0$, $\sqrt{|g|} = |\det\hb_{s}|^{-1}$ diverging. If ${*}_{g}$ had a finite limit $S$, then $I\,\mathcal{G}_{s} = |\det\hb_{s}|\,{*}_{g}$ would tend to $0\cdot S = 0$; but it tends to $I\,\mathcal{G}_{0} \neq 0$. So ${*}_{g}$ has no limit, and its bounded rescaling $|\det\hb_{s}|\,{*}_{g} \to I\,\mathcal{G}_{0}$ has rank three on the six-dimensional bivectors (the bivectors of the three-dimensional image of $\mathcal{G}_{0}$): ${*}$ on two-forms is not a nondegenerate map there, for every representative of the conformal class alike, since the scale cancels. The $\theta$-terms $\theta\int\langle\mathcal{F}\wedge\mathcal{F}\rangle_{\mathrm{int}}$ use no $\hb$, hence cross.

The check confirms in exact arithmetic the Bianchi identity and $\langle\mathcal{F}\wedge\mathcal{F}\rangle_{\mathrm{int}} = \nabla\wedge\mathcal{C}$ for polynomial $\mathfrak{su}(2)$-valued fields (bivectors of $\Cl(3,0)$, one abelian and one nonabelian sample), the formula for ${*}_{g}$ against the component star, its scale and rotation-gauge invariance and ${*}_{g}^{2} = -1$, and the divergence and rank three at a wall.

\emph{In components} (a reading): ${*}\alpha_{ab} = \tfrac12\sqrt{|g|}\,\epsilon_{abcd}\,g^{ce}g^{df}\alpha_{ef}$, $\sqrt{|g|}$ scaling by $\lambda^{-4}$ and the two inverse metrics by $\lambda^{4}$ under $g \mapsto \lambda^{-2}g$; at the wall $g^{ce}$ has a one-dimensional radical and $\sqrt{|g|}$ a half-power pole; and $d\,\tr(F\wedge F) = 2\tr(DF\wedge F) = 0$, $\tr(F\wedge F) = d\,\mathrm{CS}(A)$, with the internal algebra as matrices and the trace for $\langle\cdot\rangle_{\mathrm{int}}$ (the earlier checks compute these for $\mathfrak{gl}(2)$).
